```tex
\section{Konzeption des chemischen Analogcomputermodells zur Flächenberechnung}

\subsection{Ziel der Modellierung}

Aus der analytischen Vorbetrachtung ergibt sich für den maximalen
Flächeninhalt

\[
A_{\max}(s)
=
\frac{\pi+\sqrt{\pi^2+4}}{4}s^2.
\]

Zur Vereinfachung definieren wir die Konstante

\[
C
=
\frac{\pi+\sqrt{\pi^2+4}}{4}.
\]

Damit lautet die zu implementierende Berechnungsvorschrift

\[
A_{\max}(s)=C s^2.
\]

Numerisch gilt

\[
C \approx 1.716446108.
\]

Das chemische Analogcomputermodell soll also nicht mehr die ursprüngliche
geometrische Optimierung durchführen. Diese wurde bereits analytisch
gelöst. Stattdessen soll das Reaktionssystem die daraus resultierende
Berechnungsformel

\[
s \mapsto C s^2
\]

dynamisch realisieren.

Die Eingangsgröße ist die Schenkellänge \(s\). Diese wird durch die
Konzentration einer Eingabespezies \(S\) kodiert:

\[
[S]=s.
\]

Die Ausgabegröße ist der maximale Flächeninhalt. Dieser wird durch die
Konzentration einer Ausgabespezies \(F\) kodiert:

\[
[F]\longrightarrow A_{\max}(s).
\]

Das Ziel ist daher, ein Reaktionssystem so zu konstruieren, dass im
stationären Zustand gilt:

\[
[F]^*
=
C[S]^2.
\]

Dabei bezeichnet \([F]^*\) die Gleichgewichtskonzentration der
Ausgabespezies \(F\).

\subsection{Grundprinzip: Rechnen durch Konzentrationsdynamik}

In einem chemischen Analogcomputermodell werden Zahlenwerte durch
Konzentrationen chemischer Spezies dargestellt. Eine Berechnung wird nicht
durch eine direkte algebraische Formel ausgewertet, sondern durch die
Zeitentwicklung eines Reaktionsnetzwerks erzeugt.

Das bedeutet: Die Ausgabe ist im Allgemeinen nicht sofort zu Beginn der
Simulation korrekt. Stattdessen starten die Ausgabespezies mit
vorgegebenen Anfangskonzentrationen, zum Beispiel bei \(0\), und nähern
sich im Laufe der Modellzeit den gewünschten Zielwerten an.

Für die Flächenberechnung soll also eine Spezies \(F\) konstruiert werden,
deren Konzentration gegen

\[
C s^2
\]

konvergiert.

Die Berechnung erfolgt modular. Zuerst wird aus der Eingabe \(s\) der
Zwischenwert \(s^2\) erzeugt. Anschließend wird dieser Zwischenwert mit
der Konstanten \(C\) skaliert.

Die Rechenstruktur lautet daher:

\[
S
\longrightarrow
Q
\longrightarrow
F.
\]

Dabei steht

\[
[S]=s
\]

für die Eingabe,

\[
[Q]\longrightarrow s^2
\]

für den quadratischen Zwischenwert und

\[
[F]\longrightarrow C s^2
\]

für den maximalen Flächeninhalt.

\subsection{Warum ein dynamisches Reaktionsmodell verwendet wird}

Eine rein algebraische Modellierung wäre prinzipiell möglich. Man könnte
die Ausgabe beispielsweise direkt durch

\[
F = C S^2
\]

definieren. Für die vorliegende Aufgabenstellung wäre dies jedoch
inhaltlich unpassend, weil dann keine chemische Rechenstruktur entsteht.
Das Ziel ist nicht nur die numerische Auswertung der Formel, sondern die
Konstruktion eines chemischen Analogcomputermodells.

Deshalb wird die Berechnung als dynamisches Reaktionssystem modelliert.
Die mathematische Idee besteht darin, Spezies einzuführen, deren
Konzentrationen stabile Gleichgewichtswerte besitzen. Diese
Gleichgewichtswerte entsprechen den gewünschten Rechenergebnissen.

Allgemein kann eine gewünschte Zielbeziehung

\[
X^* = g(\text{Eingaben})
\]

durch eine Differentialgleichung der Form

\[
\frac{dX}{dt}
=
g(\text{Eingaben})-X
\]

realisiert werden. Im stationären Zustand gilt

\[
\frac{dX}{dt}=0.
\]

Daraus folgt

\[
0=g(\text{Eingaben})-X
\]

und somit

\[
X^*=g(\text{Eingaben}).
\]

Der Term \(g(\text{Eingaben})\) wirkt als Produktionsterm, während der Term
\(-X\) einen Abbau oder eine negative Rückkopplung beschreibt. Dadurch
konvergiert \(X\) gegen den gewünschten Rechenwert.

\subsection{Eingabespezies \(S\)}

Die Schenkellänge \(s\) ist in der Aufgabenstellung eine frei wählbare,
aber während einer Simulation feste Eingabegröße. Daher wird sie als
Konzentration einer Eingabespezies \(S\) kodiert:

\[
[S](0)=s.
\]

Da \(s\) nicht durch die Berechnung verbraucht werden soll, wird \(S\) als
konstante Eingabe behandelt. Das Reaktionssystem darf also von \(S\)
abhängen, soll die Konzentration von \(S\) aber nicht verändern.

Konzeptionell entspricht \(S\) damit einem Parameter oder einer
chemostatisierten Spezies. Für jede gewählte Schenkellänge \(s\) wird eine
entsprechende Konzentration von \(S\) vorgegeben. Das gleiche
Reaktionsnetzwerk kann anschließend für verschiedene Werte von \(s\)
simuliert werden.

Beispiele wären

\[
[S]=0.5,\qquad [S]=1,\qquad [S]=2,\qquad [S]=5.
\]

Die Ausgabe \(F\) sollte dann jeweils gegen

\[
C\cdot 0.5^2,\qquad
C\cdot 1^2,\qquad
C\cdot 2^2,\qquad
C\cdot 5^2
\]

konvergieren.

\subsection{Quadratmodul}

Der erste eigentliche Rechenschritt ist die Quadratbildung

\[
s \mapsto s^2.
\]

Dazu wird eine Zwischenspezies \(Q\) eingeführt. Ihre Konzentration soll im
stationären Zustand dem Quadrat der Eingabekonzentration entsprechen:

\[
[Q]^*=[S]^2.
\]

Dies kann durch die Dynamik

\[
\frac{d[Q]}{dt}
=
[S]^2-[Q]
\]

realisiert werden.

Der Produktionsterm

\[
[S]^2
\]

erzeugt \(Q\) mit einer Rate, die quadratisch von der Eingabe abhängt. Der
Abbauterm

\[
-[Q]
\]

verhindert ein unbegrenztes Anwachsen von \(Q\) und stabilisiert den
Rechenwert.

Im stationären Zustand gilt

\[
\frac{d[Q]}{dt}=0.
\]

Damit folgt

\[
0=[S]^2-[Q].
\]

Also ist

\[
[Q]^*=[S]^2.
\]

Da \([S]=s\) gilt, erhält man

\[
[Q]^*=s^2.
\]

Damit implementiert \(Q\) den benötigten Zwischenwert \(s^2\).

\subsection{Interpretation des Quadratmoduls als Reaktionssystem}

Die Differentialgleichung

\[
\frac{d[Q]}{dt}
=
[S]^2-[Q]
\]

kann als Zusammenspiel zweier Reaktionsterme interpretiert werden.

Der erste Term erzeugt \(Q\) mit einer Rate proportional zu \([S]^2\):

\[
\varnothing \longrightarrow Q.
\]

Die zugehörige Reaktionsrate ist konzeptionell

\[
v_1=[S]^2.
\]

Der zweite Term baut \(Q\) mit einer Rate proportional zu seiner eigenen
Konzentration ab:

\[
Q \longrightarrow \varnothing.
\]

Die zugehörige Reaktionsrate ist

\[
v_2=[Q].
\]

Die Nettoänderung von \(Q\) ist daher

\[
\frac{d[Q]}{dt}=v_1-v_2=[S]^2-[Q].
\]

Dieses Modul besitzt eine stabile Gleichgewichtskonzentration. Ist
\([Q]\) kleiner als \([S]^2\), dann ist

\[
[S]^2-[Q]>0,
\]

also steigt \([Q]\). Ist \([Q]\) größer als \([S]^2\), dann ist

\[
[S]^2-[Q]<0,
\]

also sinkt \([Q]\). Dadurch reguliert sich die Konzentration von \(Q\)
selbstständig auf den gewünschten Wert ein.

\subsection{Flächenmodul}

Nach der Quadratbildung muss der Zwischenwert \(s^2\) mit der Konstanten

\[
C=
\frac{\pi+\sqrt{\pi^2+4}}{4}
\]

multipliziert werden. Da \([Q]\) gegen \(s^2\) konvergiert, soll die
Ausgabespezies \(F\) im stationären Zustand die Beziehung

\[
[F]^*=C[Q]^*
\]

erfüllen.

Eine naheliegende dynamische Realisierung wäre

\[
\frac{d[F]}{dt}
=
C[Q]-[F].
\]

Im stationären Zustand gilt dann

\[
0=C[Q]-[F],
\]

also

\[
[F]^*=C[Q]^*.
\]

Da

\[
[Q]^*=s^2
\]

folgt daraus

\[
[F]^*=Cs^2=A_{\max}(s).
\]

Damit wäre die gewünschte Flächenberechnung erreicht.

\subsection{Berücksichtigung der Begrenzung der Ratenkonstanten}

Die Aufgabenstellung fordert, dass keine Ratenkonstante im verwendeten
Reaktionssystem einen Wert größer als \(1\) besitzen darf. Die Konstante
\(C\) ist jedoch

\[
C\approx1.716446108
\]

und damit größer als \(1\). Deshalb darf der Term \(C[Q]\) nicht einfach
als einzelner Reaktionsterm mit der Ratenkonstante \(C\) umgesetzt werden.

Stattdessen wird \(C\) in zwei Summanden zerlegt:

\[
C=1+c
\]

mit

\[
c=C-1.
\]

Einsetzen des Zahlenwerts ergibt

\[
c\approx0.716446108.
\]

Damit gilt

\[
C[Q]
=
(1+c)[Q]
=
[Q]+c[Q].
\]

Die Produktionsrate für \(F\) wird also nicht als ein einzelner Term mit
der Ratenkonstante \(C>1\) modelliert, sondern als Summe zweier erlaubter
Produktionsterme:

\[
[Q]
\]

und

\[
c[Q].
\]

Beide zugehörigen Ratenkonstanten erfüllen die Bedingung

\[
1\leq 1
\]

und

\[
c\approx0.716446108<1.
\]

Die Dynamik des Flächenmoduls lautet somit

\[
\frac{d[F]}{dt}
=
[Q]+c[Q]-[F].
\]

Dies ist äquivalent zu

\[
\frac{d[F]}{dt}
=
(1+c)[Q]-[F]
\]

und wegen \(1+c=C\) zu

\[
\frac{d[F]}{dt}
=
C[Q]-[F].
\]

Die Zerlegung verändert also nicht die mathematische Zielbeziehung,
stellt aber sicher, dass keine einzelne Ratenkonstante größer als \(1\)
verwendet wird.

\subsection{Interpretation des Flächenmoduls als Reaktionssystem}

Das Flächenmodul kann als drei Reaktionsterme interpretiert werden.

Zunächst wird \(F\) mit einer Rate proportional zu \([Q]\) produziert:

\[
\varnothing \longrightarrow F,
\]

mit

\[
v_3=[Q].
\]

Zusätzlich wird \(F\) über einen zweiten parallelen Produktionskanal
erzeugt:

\[
\varnothing \longrightarrow F,
\]

mit

\[
v_4=c[Q].
\]

Schließlich wird \(F\) mit einer Rate proportional zu seiner eigenen
Konzentration abgebaut:

\[
F\longrightarrow\varnothing,
\]

mit

\[
v_5=[F].
\]

Die Nettoänderung von \(F\) ergibt sich damit zu

\[
\frac{d[F]}{dt}
=
v_3+v_4-v_5.
\]

Also

\[
\frac{d[F]}{dt}
=
[Q]+c[Q]-[F].
\]

Im stationären Zustand gilt

\[
0=[Q]+c[Q]-[F].
\]

Daher ist

\[
[F]^*=(1+c)[Q]^*.
\]

Da

\[
1+c=C
\]

folgt

\[
[F]^*=C[Q]^*.
\]

Mit

\[
[Q]^*=s^2
\]

erhält man schließlich

\[
[F]^*=Cs^2.
\]

Also gilt

\[
[F]^*
=
\frac{\pi+\sqrt{\pi^2+4}}{4}s^2
=
A_{\max}(s).
\]

Damit ist die Berechnung des maximalen Flächeninhalts durch das
Reaktionssystem realisiert.

\subsection{Gesamtdynamik des Flächenmodells}

Das vollständige Modell für die Flächenberechnung besteht aus den beiden
gekoppelten Differentialgleichungen

\[
\frac{d[Q]}{dt}
=
[S]^2-[Q]
\]

und

\[
\frac{d[F]}{dt}
=
[Q]+c[Q]-[F].
\]

Äquivalent kann man schreiben:

\[
\frac{d[Q]}{dt}
=
[S]^2-[Q],
\]

\[
\frac{d[F]}{dt}
=
C[Q]-[F],
\]

wobei bei der tatsächlichen Reaktionsmodellierung der Faktor \(C\) durch
die Zerlegung

\[
C=1+c
\]

realisiert wird.

Die Eingabespezies \(S\) bleibt konstant:

\[
\frac{d[S]}{dt}=0.
\]

Damit ergibt sich das Gesamtsystem

\[
\boxed{
\frac{d[S]}{dt}=0
}
\]

\[
\boxed{
\frac{d[Q]}{dt}=[S]^2-[Q]
}
\]

\[
\boxed{
\frac{d[F]}{dt}=[Q]+c[Q]-[F]
}
\]

mit

\[
c
=
\frac{\pi+\sqrt{\pi^2+4}}{4}-1.
\]

Im stationären Zustand gilt

\[
[Q]^*=[S]^2
\]

und

\[
[F]^*=C[Q]^*.
\]

Daraus folgt

\[
[F]^*=C[S]^2.
\]

Da

\[
[S]=s
\]

gilt, erhält man

\[
[F]^*
=
C s^2
=
A_{\max}(s).
\]

\subsection{Warum eine Kaskadenstruktur verwendet wird}

Die Modellierung erfolgt bewusst als Kaskade

\[
S
\longrightarrow
Q
\longrightarrow
F.
\]

Diese Struktur hat mehrere Vorteile.

Erstens trennt sie die beiden Rechenschritte klar voneinander. Das
Quadratmodul berechnet zunächst den Zwischenwert \(s^2\). Das
Flächenmodul skaliert diesen Zwischenwert anschließend mit der Konstanten
\(C\).

Zweitens ist das Modell dadurch leichter überprüfbar. Bei einer Simulation
kann man separat kontrollieren, ob

\[
[Q]\longrightarrow s^2
\]

und danach

\[
[F]\longrightarrow C s^2
\]

gilt.

Drittens ist die Kaskade gut erweiterbar. Falls später weitere
Berechnungsvorschriften oder zusätzliche Ausgaben ergänzt werden sollen,
kann der Zwischenwert \(Q\) weiterverwendet werden.

Viertens ist die spätere Analyse der chemischen Rechenzeit einfacher.
Da \(F\) von \(Q\) abhängt, ist zu erwarten, dass \(F\) etwas langsamer
gegen seinen Zielwert konvergiert als \(Q\). Diese Verzögerung ist ein
natürlicher Effekt der Kaskadenstruktur und kann in den
Simulationsstudien quantitativ untersucht werden.

\subsection{Rolle der Abbaureaktionen}

Die Abbaureaktionen sind ein zentraler Bestandteil des Modells. Ohne
Abbau würden die Produkte \(Q\) und \(F\) unbegrenzt anwachsen.

Für das Quadratmodul würde ohne Abbau beispielsweise gelten:

\[
\frac{d[Q]}{dt}=[S]^2.
\]

Bei konstantem \(S\) hätte dies die Lösung

\[
[Q](t)=[Q](0)+[S]^2t.
\]

Die Konzentration von \(Q\) würde also linear mit der Zeit wachsen und
keinen stabilen Rechenwert darstellen.

Erst durch den zusätzlichen Term

\[
-[Q]
\]

entsteht die stabile Dynamik

\[
\frac{d[Q]}{dt}=[S]^2-[Q].
\]

Analog verhindert der Term

\[
-[F]
\]

im Flächenmodul, dass \(F\) unbegrenzt anwächst. Stattdessen wird \(F\)
auf den Zielwert

\[
C[Q]^*
\]

eingeregelt.

Die Abbaureaktionen wirken daher als negative Rückkopplung. Sie sorgen
dafür, dass die Rechenspezies stabile Gleichgewichtskonzentrationen
besitzen.

\subsection{Deterministische Modellierung}

Für die erste Modellvariante wird eine deterministische Beschreibung
verwendet. Dabei werden die Konzentrationen als kontinuierliche Größen
aufgefasst, deren Zeitentwicklung durch gewöhnliche Differentialgleichungen
beschrieben wird.

Diese Wahl ist für das vorliegende Modell sinnvoll, weil die Aufgabe eine
analoge Berechnung über Konzentrationen betrachtet. Die relevanten Größen
sind keine einzelnen diskreten Moleküle, sondern kontinuierliche
Konzentrationswerte, die gegen mathematische Zielwerte konvergieren.

Eine stochastische Simulation wäre prinzipiell ebenfalls denkbar. Sie
würde jedoch zufällige Fluktuationen einführen und die Auswertung der
Rechengenauigkeit erschweren. Für die grundlegende Implementierung der
Berechnungsvorschrift

\[
A_{\max}(s)=Cs^2
\]

ist daher zunächst das deterministische ODE-Modell die geeignetere Wahl.

\subsection{Zusammenfassung des Flächenmodells}

Das chemische Analogcomputermodell zur Flächenberechnung verwendet die
Eingabespezies \(S\), die Zwischenspezies \(Q\) und die Ausgabespezies
\(F\).

Die Konzentration von \(S\) kodiert die Schenkellänge:

\[
[S]=s.
\]

Die Spezies \(Q\) berechnet das Quadrat der Eingabe:

\[
[Q]^*=[S]^2=s^2.
\]

Die Spezies \(F\) berechnet daraus den maximalen Flächeninhalt:

\[
[F]^*
=
C[Q]^*
=
Cs^2.
\]

Mit

\[
C=
\frac{\pi+\sqrt{\pi^2+4}}{4}
\]

ergibt sich

\[
[F]^*
=
\frac{\pi+\sqrt{\pi^2+4}}{4}s^2
=
A_{\max}(s).
\]

Da

\[
C\approx1.716446108>1
\]

ist, wird der Faktor \(C\) nicht durch eine einzelne Ratenkonstante
realisiert. Stattdessen wird

\[
C=1+c
\]

mit

\[
c\approx0.716446108
\]

verwendet. Dadurch bleiben alle Ratenkonstanten höchstens \(1\), wie es
die Aufgabenstellung fordert.

Die zentrale Modellstruktur lautet damit:

\[
S
\longrightarrow
Q
\longrightarrow
F,
\]

wobei \(S\) konstant gehalten wird, \(Q\) gegen \(s^2\) konvergiert und
\(F\) gegen den maximalen Flächeninhalt \(A_{\max}(s)\) konvergiert.
```
